\section{Thiết kế thực nghiệm}

\subsection{Tập dữ liệu và giao thức chia tập}

Thực nghiệm sử dụng tập dữ liệu tùy chỉnh đã được chuẩn bị trước và đồng bộ lên Amazon S3. Dữ liệu đầu vào đã bao gồm các biến thể điều kiện môi trường cần thiết; benchmark AWS chỉ đọc dữ liệu và chạy inference, không sinh augmentation lại trong quá trình đo.

Sau bước sinh cặp và lọc các mẫu hợp lệ qua pipeline nhận diện, tập đánh giá gồm 248 cặp xác minh. Tập này được chia ngẫu nhiên nhưng cố định bằng seed 42:

\begin{itemize}
    \item \textbf{Calibration split}: 124 cặp, dùng để hiệu chỉnh ngưỡng theo FAR budget.
    \item \textbf{Test split}: 124 cặp, dùng để báo cáo kết quả cuối cùng.
\end{itemize}

Việc tách calibration/test giúp tránh đánh giá ngưỡng ngay trên dữ liệu đã dùng để chọn ngưỡng. Tuy nhiên, quy mô 248 cặp vẫn nhỏ, vì vậy các kết quả cần được xem là bằng chứng thực nghiệm ban đầu thay vì kết luận tổng quát.

\subsection{Mô hình nhận diện}

Mô hình nhận diện dùng InsightFace \textit{buffalo\_sc}. Trong log chạy thực nghiệm, mô hình recognition được nạp từ file \textit{w600k\_mbf.onnx}, tức backbone MBF/WebFace600K mang phong cách MobileFaceNet. Đây là lựa chọn phù hợp hơn với mục tiêu edge so với các gói InsightFace lớn hơn. Detector chạy với kích thước đầu vào $320 \times 320$, sử dụng CPUExecutionProvider của ONNX Runtime.

Cấu hình chính:

\begin{itemize}
    \item Requested model: \texttt{mobilefacenet}
    \item Resolved InsightFace pack: \texttt{buffalo\_sc}
    \item Recognition backend: \texttt{w600k\_mbf.onnx}
    \item Detector size: $320 \times 320$
    \item Embedding backend: InsightFace/ONNX Runtime CPU
\end{itemize}

\subsection{Các FAR budget}

Thực nghiệm risk-constrained đánh giá bốn mức ngân sách rủi ro:

\begin{equation}
    \alpha \in \{0{,}01,\;0{,}02,\;0{,}03,\;0{,}05\}.
\end{equation}

Với mỗi $\alpha$, \method{M5} và \method{M6} hiệu chỉnh ngưỡng riêng theo condition bin trên tập calibration. Với \method{M6}, margin defer là $\delta=0{,}03$.

\subsection{Mô phỏng thiết bị biên trên AWS}

Benchmark được chạy trên AWS EC2 ARM và dùng Docker để giới hạn CPU/RAM. Cách làm này không giống hoàn toàn Raspberry Pi hoặc Jetson thật, nhưng tạo ra môi trường có thể tái lập để so sánh các phương pháp dưới cùng ràng buộc tài nguyên.

\begin{table}[!htbp]
\centering
\caption{Các cấu hình edge-constrained được dùng trong benchmark AWS}
\label{tab:aws_profiles}
\begin{tabular}{lcc}
\toprule
\textbf{Profile} & \textbf{CPU limit} & \textbf{Memory limit} \\
\midrule
\texttt{edge\_1cpu\_512mb} & 1 CPU & 512 MB \\
\texttt{edge\_1cpu\_1gb}   & 1 CPU & 1 GB \\
\texttt{edge\_2cpu\_2gb}   & 2 CPU & 2 GB \\
\bottomrule
\end{tabular}
\end{table}

Các profile này được chạy bằng cùng một script benchmark để đo cả hiệu năng nhận diện và chi phí tài nguyên. Kết quả được lưu ở hai nơi:

\begin{itemize}
    \item Local: \texttt{D:\textbackslash face\_recognition\textbackslash aws\_edge\_benchmark\_results}
    \item S3: \texttt{s3://face-edge-sim-380269593136-ap-southeast-1/aws-edge-benchmark/latest/}
\end{itemize}

\subsection{File kết quả}

Các bảng trong Mục~\ref{sec:results} được lấy từ các file CSV sau:

\begin{itemize}
    \item \texttt{combined\_edge\_summary.csv}: tổng hợp theo method/profile.
    \item \texttt{combined\_far\_budget\_summary.csv}: tổng hợp theo FAR budget và condition bin.
    \item \texttt{combined\_calibration\_thresholds.csv}: ngưỡng accept/reject sau calibration.
\end{itemize}

Trong báo cáo, kết quả nhận diện chính được trình bày trên profile \texttt{edge\_1cpu\_512mb} vì đây là cấu hình chặt nhất. Bảng profile riêng được dùng để so sánh độ trễ và RAM trên ba cấu hình AWS.
